\section{引言}
\sourcepage{443}
如前所述，流形是数学中最熟悉、最重要的一类几何对象。数学家经过相当长的时间，才给出流形概念的严格定义。最终，Hermann Weyl 在其关于 Riemann 曲面的著作中明确了这一概念。\footnote{\textit{Die Idee der Riemannschen Fläche}，Teubner，Leipzig--Berlin，1913。第三版的英译本见参考文献 [Wey64]。}

关于流形所证明的定理，往往不仅对拓扑学有用，而且对微分几何及其他数学分支也有用。对偶定理正是这样的例子。

对于任意拓扑空间，同调群与上同调群之间已有一种关系，即上同调的泛系数定理所表达的关系。对流形而言，还存在第二种关系，这就是著名的 Poincaré 对偶定理。如果把泛系数定理看作表达上同调与同调之间的一种代数对偶性，那么 Poincaré 对偶性从根本上说可以看作一种几何性质。这种对偶性并不对任意空间成立，而是特别依赖于流形所具有的性质。

这两类几何与代数的对偶性之间存在相互作用。这种相互作用有一些有趣的应用；我们将把它用于计算流形的上同调环。

还有其他对偶定理，分别以 Lefschetz、Alexander 和 Pontryagin 的名字命名。我们将在适当时候研究它们。这些定理都涉及流形的代数拓扑性质。

在以下讨论中，假定读者熟悉局部同调群（第 4.8 节）。
